\documentclass[10pt,a4paper]{amsart}
\usepackage{amsmath,amssymb,amsthm,graphicx}
% Retained preamble for the Stacks Project chapter {CH} (shared project preamble).
% The chapter file pinned-original-source.tex only carries \input{{preamble}} on its first line,
% so no preamble text lives inside the chapter itself.  The Stacks Project uses one single shared
% preamble for every chapter: the snapshot stages it three times (sources/stacks-chow/preamble.tex,
% sources/stacks-etale-cohomology/preamble.tex and focused-seeds/AGP/sources/stacks-snapshot-20261011/
% preamble.tex) and those copies are identical in size.  The block below is copied mechanically,
% line for line, from sources/stacks-chow/preamble.tex lines 14-18 and 154-217: the Xy-pic setup and
% the native theorem declarations together with the mathematical macros.  Nothing was added, removed
% or redefined; the wrapper class and the amsmath/amssymb/amsthm/graphicx base packages are supplied
% by the assembler.
% For commutative diagrams we use Xy-pic
% Chapter: stacks-cohomology.tex -- Cohomology of Sheaves.
\usepackage[all]{xy}

% We use 2cell for 2-commutative diagrams.
\xyoption{2cell}

% Theorem environments.
%
\theoremstyle{plain}
\newtheorem{theorem}[subsection]{Theorem}
\newtheorem{proposition}[subsection]{Proposition}
\newtheorem{lemma}[subsection]{Lemma}

\theoremstyle{definition}
\newtheorem{definition}[subsection]{Definition}
\newtheorem{example}[subsection]{Example}
\newtheorem{exercise}[subsection]{Exercise}
\newtheorem{situation}[subsection]{Situation}

\theoremstyle{remark}
\newtheorem{remark}[subsection]{Remark}
\newtheorem{remarks}[subsection]{Remarks}

\numberwithin{equation}{subsection}

% Macros
%
\def\lim{\mathop{\mathrm{lim}}\nolimits}
\def\colim{\mathop{\mathrm{colim}}\nolimits}
\def\Spec{\mathop{\mathrm{Spec}}}
\def\Hom{\mathop{\mathrm{Hom}}\nolimits}
\def\Ext{\mathop{\mathrm{Ext}}\nolimits}
\def\SheafHom{\mathop{\mathcal{H}\!\mathit{om}}\nolimits}
\def\SheafExt{\mathop{\mathcal{E}\!\mathit{xt}}\nolimits}
\def\Sch{\mathit{Sch}}
\def\Mor{\mathop{\mathrm{Mor}}\nolimits}
\def\Ob{\mathop{\mathrm{Ob}}\nolimits}
\def\Sh{\mathop{\mathit{Sh}}\nolimits}
\def\NL{\mathop{N\!L}\nolimits}
\def\CH{\mathop{\mathrm{CH}}\nolimits}
\def\proetale{{pro\text{-}\acute{e}tale}}
\def\etale{{\acute{e}tale}}
\def\QCoh{\mathit{QCoh}}
\def\Ker{\mathop{\mathrm{Ker}}}
\def\Im{\mathop{\mathrm{Im}}}
\def\Coker{\mathop{\mathrm{Coker}}}
\def\Coim{\mathop{\mathrm{Coim}}}

% Boxtimes
%
\DeclareMathSymbol{\boxtimes}{\mathbin}{AMSa}{"02}

%
% Macros for moduli stacks/spaces
%
\def\QCohstack{\mathcal{QC}\!\mathit{oh}}
\def\Cohstack{\mathcal{C}\!\mathit{oh}}
\def\Spacesstack{\mathcal{S}\!\mathit{paces}}
\def\Quotfunctor{\mathrm{Quot}}
\def\Hilbfunctor{\mathrm{Hilb}}
\def\Curvesstack{\mathcal{C}\!\mathit{urves}}
\def\Polarizedstack{\mathcal{P}\!\mathit{olarized}}
\def\Complexesstack{\mathcal{C}\!\mathit{omplexes}}
% \Pic is the operator that assigns to X its picard group, usage \Pic(X)
% \Picardstack_{X/B} denotes the Picard stack of X over B
% \Picardfunctor_{X/B} denotes the Picard functor of X over B
\def\Pic{\mathop{\mathrm{Pic}}\nolimits}
\def\Picardstack{\mathcal{P}\!\mathit{ic}}
\def\Picardfunctor{\mathrm{Pic}}
\def\Deformationcategory{\mathcal{D}\!\mathit{ef}}

% --- appended 2026-10-12 (Chief): the project own definitions of the reference/slogan/history
% display environments, copied mechanically from the shared Stacks preamble lines 8-12
% (\usepackage{verbatim} plus the three \newenvironment lines, which the chapter preamble keeps
% verbatim from the project preamble). Nothing added, removed or redefined: these environments
% carry no mathematical content and the project itself typesets them as comments. Without these
% three lines a counted statement whose own text uses them fails with Environment slogan undefined.
\usepackage{verbatim}
\newenvironment{reference}{\comment}{\endcomment}
\newenvironment{slogan}{\comment}{\endcomment}
\newenvironment{history}{\comment}{\endcomment}
\usepackage{cleveref}
\crefname{originalreference}{original reference}{original references}
\Crefname{originalreference}{Original reference}{Original references}
\makeatletter
\newlabel{situation-locally-given}{{0D67}{source}{}{sourcepointer-1}{}}
\newlabel{situation-locally-given@cref}{{[originalreference][0][]0D67}{[1][0][]source}{}{}{}}
\makeatother
\crefname{originallemma}{original lemma}{original lemmas}
\Crefname{originallemma}{Original lemma}{Original lemmas}
\makeatletter
\newlabel{lemma-uniqueness}{{0D68}{source}{}{sourcepointer-2}{}}
\newlabel{lemma-uniqueness@cref}{{[originallemma][0][]0D68}{[1][0][]source}{}{}{}}
\makeatother
\crefname{originalremark}{original remark}{original remarks}
\Crefname{originalremark}{Original remark}{Original remarks}
\makeatletter
\newlabel{remark-uniqueness}{{0D69}{source}{}{sourcepointer-3}{}}
\newlabel{remark-uniqueness@cref}{{[originalremark][0][]0D69}{[1][0][]source}{}{}{}}
\makeatother
\crefname{originallemma}{original lemma}{original lemmas}
\Crefname{originallemma}{Original lemma}{Original lemmas}
\makeatletter
\newlabel{lemma-glue}{{08DG}{source}{}{sourcepointer-4}{}}
\newlabel{lemma-glue@cref}{{[originallemma][0][]08DG}{[1][0][]source}{}{}{}}
\makeatother
\title{Gluing solutions in the finite case: existence and uniqueness up to unique isomorphism}
\author{The Stacks Project Authors}
\date{stacks-project:cohomology.tex:local-snapshot-20261011; GNU FDL 1.2 or later}
\begin{document}
\maketitle
\noindent\textit{Editorial note.} All delivered mathematical text is the author's own, copied verbatim from the pinned Stacks Project chapter file cohomology.tex; no statement or proof step was rewritten, reordered or removed.\par\medskip
\noindent\textit{Allowed standard prerequisites (editorial).}\begin{itemize}
\item delivered companion: gluing lemma. delivered companion lemma-glue (S11-csh-glue). Produces the glued object K with compatible isomorphisms in the n=2 case.
\item Same-chapter uniqueness lemma and remark. Cohomology, Lemma cohomology-lemma-uniqueness; Remark cohomology-remark-uniqueness. Uniqueness up to unique isomorphism; compatibility property pinning the identifications.
\item Standing situation (statement level). Cohomology, Situation cohomology-situation-locally-given. Locally given system (K\_U, rho) with gluing isomorphisms on a basis.
\end{itemize}\medskip
\ifdefined\hypertarget\hypertarget{sourcepointer-1}{}\fi
\noindent\textit{Source reference.} Original reference 0D67: sources/stacks-chow/tags.txt, source lines 16512--16512. Cross-chapter standard prerequisite cited by the author inside the counted span; it is used as a pointer and its own proof is not part of this item.\par\medskip
\ifdefined\hypertarget\hypertarget{sourcepointer-2}{}\fi
\noindent\textit{Source reference.} Original lemma 0D68: sources/stacks-chow/tags.txt, source lines 16513--16513. Cross-chapter standard prerequisite cited by the author inside the counted span; it is used as a pointer and its own proof is not part of this item.\par\medskip
\ifdefined\hypertarget\hypertarget{sourcepointer-3}{}\fi
\noindent\textit{Source reference.} Original remark 0D69: sources/stacks-chow/tags.txt, source lines 16514--16514. Cross-chapter standard prerequisite cited by the author inside the counted span; it is used as a pointer and its own proof is not part of this item.\par\medskip
\ifdefined\hypertarget\hypertarget{sourcepointer-4}{}\fi
\noindent\textit{Source reference.} Original lemma 08DG: sources/stacks-chow/tags.txt, source lines 10636--10636. Cross-chapter standard prerequisite cited by the author inside the counted span; it is used as a pointer and its own proof is not part of this item.\par\medskip
\begin{lemma}
\label{lemma-solution-in-finite-case}
In Situation \ref{situation-locally-given} assume
\begin{enumerate}
\item $X = U_1 \cup \ldots \cup U_n$ with $U_i \in \mathcal{B}$,
\item for $U, V \in \mathcal{B}$ we have
$U \cap V = \bigcup_{W \in \mathcal{B}, W \subset U \cap V} W$,
\item for any $U \in \mathcal{B}$ we have $\Ext^i(K_U, K_U) = 0$
for $i < 0$.
\end{enumerate}
Then a solution exists and is unique up to unique isomorphism.
\end{lemma}

\begin{proof}
Uniqueness was seen in Lemma \ref{lemma-uniqueness}. We may prove the lemma
by induction on $n$. The case $n = 1$ is immediate.

\medskip\noindent
The case $n = 2$.
Consider the isomorphism
$\rho_{U_1, U_2} : K_{U_1}|_{U_1 \cap U_2} \to K_{U_2}|_{U_1 \cap U_2}$
constructed in Remark \ref{remark-uniqueness}.
By Lemma \ref{lemma-glue} we obtain an object $K$ in $D(\mathcal{O}_X)$
and isomorphisms $\rho_{U_1} : K|_{U_1} \to K_{U_1}$ and
$\rho_{U_2} : K|_{U_2} \to K_{U_2}$ compatible with $\rho_{U_1, U_2}$.
Take $U \in \mathcal{B}$. We will construct an isomorphism
$\rho_U : K|_U \to K_U$ and we will leave it to the reader to verify
that $(K, \rho_U)$ is a solution. Consider the set $\mathcal{B}'$
of elements of $\mathcal{B}$ contained in either $U \cap U_1$ or contained in
$U \cap U_2$. Then $(K_U, \rho^U_{U'})$ is a solution for the system
$(\{K_{U'}\}_{U' \in \mathcal{B}'},
\{\rho_{V'}^{U'}\}_{V' \subset U'\text{ with }U', V' \in \mathcal{B}'})$
on the ringed space $U$.
We claim that $(K|_U, \tau_{U'})$ is another solution where
$\tau_{U'}$ for $U' \in \mathcal{B}'$ is chosen as follows:
if $U' \subset U_1$ then we take the composition
$$
K|_{U'} \xrightarrow{\rho_{U_1}|_{U'}}
K_{U_1}|_{U'} \xrightarrow{\rho^{U_1}_{U'}}
K_{U'}
$$
and if $U' \subset U_2$ then we take the composition
$$
K|_{U'} \xrightarrow{\rho_{U_2}|_{U'}}
K_{U_2}|_{U'} \xrightarrow{\rho^{U_2}_{U'}}
K_{U'}.
$$
To verify this is a solution use the property of the map $\rho_{U_1, U_2}$
described in Remark \ref{remark-uniqueness} and the compatibility of
$\rho_{U_1}$ and $\rho_{U_2}$ with $\rho_{U_1, U_2}$. Having said this
we apply Lemma \ref{lemma-uniqueness} to see that we obtain a unique
isomorphism $K|_{U'} \to K_{U'}$ compatible with the maps $\tau_{U'}$ and
$\rho^U_{U'}$ for $U' \in \mathcal{B}'$.

\medskip\noindent
The case $n > 2$. Consider the open subspace
$X' = U_1 \cup \ldots \cup U_{n - 1}$ and let $\mathcal{B}'$ be the set of
elements of $\mathcal{B}$ contained in $X'$. Then we find a system
$(\{K_U\}_{U \in \mathcal{B}'}, \{\rho_V^U\}_{U, V \in \mathcal{B}'})$
on the ringed space $X'$ to which we may apply our induction hypothesis.
We find a solution $(K_{X'}, \rho^{X'}_U)$.
Then we can consider the collection
$\mathcal{B}^* = \mathcal{B} \cup \{X'\}$ of opens of $X$ and we see that
we obtain a system
$(\{K_U\}_{U \in \mathcal{B}^*},
\{\rho_V^U\}_{V \subset U\text{ with }U, V \in \mathcal{B}^*})$.
Note that this new system also satisfies condition (3)
by Lemma \ref{lemma-uniqueness} applied to the solution $K_{X'}$.
For this system we have $X = X' \cup U_n$.
This reduces us to the case $n = 2$ we worked out above.
\end{proof}

\end{document}
